\documentclass[a4paper,10.5pt]{article}
\usepackage[utf8]{inputenc}
\usepackage[T1]{fontenc}
\usepackage[french]{babel}
\usepackage[left=1cm,right=1cm,top=.5cm,bottom=0cm]{geometry}
\usepackage{enumitem}
\usepackage{hyperref}
\usepackage{xcolor}
\usepackage{titlesec}
\usepackage{fontawesome5}
\usepackage{lmodern}
\usepackage[normalem]{ulem}

% --- Couleurs Accenture (violet corporate) ---
\definecolor{primary}{RGB}{117, 59, 189}
\definecolor{secondary}{RGB}{80, 80, 80}

% --- Config Liens ---
\hypersetup{colorlinks=true, linkcolor=primary, filecolor=primary, urlcolor=primary}

% --- Titres ---
\titleformat{\section}{\large\bfseries\color{primary}\uppercase}{}{0em}{}[\titlerule]
\titlespacing{\section}{0pt}{8pt}{5pt}

% --- Commandes ---
\newcommand{\resumeItem}[1]{\item\small{#1}}

\begin{document}

% --- EN-TÊTE ---
\begin{center}
    {\Large \textbf{Loïc Aron MBASSI EWOLO}} \\ \vspace{4pt}
    \textbf{\large Stage Développeur IA Générative} \\
    \textit{\small Disponible dès \textbf{avril 2026} pour 4 à 6 mois} \\ \vspace{4pt}
    \small
    \faMapMarker\ Cergy, France (Mobile France entière) \quad
    \faPhone\ (+33) 7 59 52 33 98 \quad
    \faEnvelope\ \href{mailto:loicaron.mbassiewolo@ensea.fr}{loicaron.mbassiewolo@ensea.fr} \\ \vspace{2pt}
    \faLinkedin\ \href{https://www.linkedin.com/in/loic-mbassi/}{linkedin.com/in/loic-mbassi} \quad
    \faGithub\ \href{https://github.com/Nameless0l}{github.com/Nameless0l} \quad
    \faGlobe\ \href{https://mbassiloic.tech}{mbassiloic.tech}
\end{center}

\vspace{-6pt}

% --- PROFIL ---
\noindent\small{Élève-ingénieur en double diplôme ENSEA / Polytechnique Yaoundé, spécialisé en \textbf{IA et Machine Learning}. Expérience concrète en \textbf{IA générative} (RAG, LangChain, prompting), \textbf{NLP} et systèmes multi-agents. Sensibilisé aux enjeux de \textbf{Responsible AI} et d'éthique. Je souhaite rejoindre \textbf{Accenture} pour contribuer à la conception d'assistants intelligents innovants.}

% --- FORMATION ---
\section{Formation}
\begin{itemize}[leftmargin=*, label=\tiny$\bullet$, nosep]
    \item \textbf{ENSEA, École Nationale Supérieure de l'Électronique et de ses Applications} \hfill \textit{2025 -- 2027} \\
    \small{Double diplôme d'Ingénieur : Systèmes Embarqués, Signal, \textbf{IA/Deep Learning}.}
    \item \textbf{École Polytechnique de Yaoundé (1\textsuperscript{ère} école d'ingénieur CEMAC)} \hfill \textit{2021 -- 2025} \\
    \small{Diplôme d'Ingénieur de Conception (Master 2) : \textbf{Mathématiques Appliquées}, Génie Logiciel, Algorithmique.}
\end{itemize}

% --- COMPÉTENCES ---
\section{Compétences Techniques}
\begin{itemize}[leftmargin=*, nosep]
    \item \small{\textbf{IA Générative \& LLM :} RAG, \textbf{LangChain}, Google ADK, Prompting, Fine-tuning, Gemini, GPT.}
    \item \small{\textbf{Machine Learning :} \textbf{PyTorch}, TensorFlow/Keras, Scikit-learn, NLP (BERT, Transformers), Computer Vision.}
    \item \small{\textbf{Langages \& Outils :} \textbf{Python}, SQL, FastAPI, Docker, Git, Streamlit, Jupyter.}
    \item \small{\textbf{Méthodologie :} Évaluation de modèles, Tests de pertinence, Responsible AI, Documentation technique.}
    \item \small{\textbf{Certifications :} Supervised Machine Learning (DeepLearning.AI/Stanford), Python for Data Science (IBM).}
\end{itemize}

% --- PROJETS IA ---
\section{Projets IA \& Machine Learning}
\begin{itemize}[leftmargin=*, label={-}]

\item \textbf{Système Multi-Agents Agricole (RAG \& LLM)} \hfill \textit{2024 -- 2025} \\
\textit{Projet Recherche} \hfill \textit{Google ADK, RAG, Gemini, LangChain, FastAPI, Docker}
\vspace{-2pt}
    \begin{itemize}[leftmargin=0.4cm, nosep, label=\tiny$\bullet$]
        \item \small{Conception d'un système de \textbf{4 agents IA collaboratifs} pour recommandations agricoles.}
        \item \small{Implémentation \textbf{RAG} pour enrichissement contextuel depuis base de connaissances.}
        \item \small{Orchestration via Google ADK, intégration APIs (météo, marchés), \textbf{prompting avancé}.}
        \item \small{Évaluation de la pertinence des réponses générées et gestion des conflits inter-agents.}
    \end{itemize}
\vspace{3pt}

\item \textbf{Laravel API Generator (Intégration LLM)} \hfill \textit{2024 -- Présent} \\
\textit{Projet Open Source (+30 étoiles GitHub)} \hfill \textit{PHP, Python, LLM, Parsing XML}
\vspace{-2pt}
    \begin{itemize}[leftmargin=0.4cm, nosep, label=\tiny$\bullet$]
        \item \small{Outil de génération automatique de code backend depuis diagrammes UML/XML.}
        \item \small{Intégration \textbf{LLM} pour assistance intelligente et génération de code contextuel.}
        \item \small{Déployé sur Packagist avec +20 téléchargements.}
    \end{itemize}
\vspace{3pt}

\item \textbf{Women Safe : Détection de Contenus Sexistes (NLP)} \hfill \textit{2023} \\
\textit{Compétition MLPC} \hfill \textit{BERT, GPT-2, Transformers, Python}
\vspace{-2pt}
    \begin{itemize}[leftmargin=0.4cm, nosep, label=\tiny$\bullet$]
        \item \small{Modèle NLP de détection de misogynie dans contenus textuels (réseaux sociaux, offres d'emploi).}
        \item \small{Étude des \textbf{biais éthiques} et approche Responsible AI.}
    \end{itemize}
\vspace{3pt}

\item \textbf{YembaAudio : Préservation Langues Locales (Low-resource NLP)} \hfill \textit{2025} \\
\textit{Projet Personnel} \hfill \textit{NLP, Speech Recognition, TensorFlow, Next.js}
\vspace{-2pt}
    \begin{itemize}[leftmargin=0.4cm, nosep, label=\tiny$\bullet$]
        \item \small{Application d'apprentissage de la langue Yemba avec reconnaissance vocale.}
        \item \small{Utilisation de modèles pré-entraînés (HuggingFace) adaptés aux langues peu dotées.}
    \end{itemize}

\end{itemize}

% --- EXPÉRIENCE PROFESSIONNELLE ---
\section{Expérience en Recherche IA}
\begin{itemize}[leftmargin=*, label={-}]

\item \textbf{Assistant de Recherche : Interprétabilité des Modèles} \hfill \textit{Juin 2025 -- Sept. 2025} \\
\textit{MILA (Institut Québécois d'IA), Remote} \hfill \textit{Python, PyTorch, TensorFlow}
\vspace{-2pt}
    \begin{itemize}[leftmargin=0.4cm, nosep, label=\tiny$\bullet$]
        \item \small{Étude du phénomène "Grokking" (généralisation tardive des réseaux de neurones).}
        \item \small{Reproduction expériences SOTA, analyse mathématique de convergence.}
        \item \small{Collaboration avec chercheurs internationaux, rigueur scientifique.}
    \end{itemize}
\vspace{3pt}

\item \textbf{Projet UROP : Analyse ECG par Computer Vision} \hfill \textit{2024} \\
\textit{Collaboration Google DeepMind} \hfill \textit{TensorFlow/Keras, OpenCV, Signal Processing}
\vspace{-2pt}
    \begin{itemize}[leftmargin=0.4cm, nosep, label=\tiny$\bullet$]
        \item \small{Modèle de détection d'anomalies cardiaques à partir d'images d'électrocardiogrammes.}
        \item \small{Extraction de features, \textbf{évaluation des performances} du modèle.}
    \end{itemize}

\end{itemize}

% --- DISTINCTIONS ---
\section{Distinctions \& Leadership}
\begin{itemize}[leftmargin=*, label=\tiny$\bullet$, nosep]
    \item \textbf{1\textsuperscript{er} Prix} IndabaX Africa 2025 (IA Santé) | \textbf{1\textsuperscript{er}} CITS National (Smart City, 75 équipes).
    \item \textbf{Lead Cellule IA} : Polytechnique Yaoundé, +50\% membres, collaboration Google DeepMind.
    \item \textbf{Animation Workshops} : Fine-tuning LLM, Machine Learning, Architecture logicielle.
    \item \textbf{Langues :} Français (natif), Anglais (professionnel/technique).
\end{itemize}

\end{document}
